\documentclass[11pt]{article}

\usepackage[utf8]{inputenc}
\usepackage[T1]{fontenc}
\usepackage{lmodern}
\usepackage{amsmath,amssymb,amsthm}
\usepackage[authoryear]{natbib}
\usepackage{hyperref}
\usepackage{geometry}
\usepackage{booktabs}
\usepackage{array}
\usepackage{parskip}
\usepackage{enumitem}
\usepackage{microtype}
\usepackage{float}
\usepackage{listings}
\usepackage{xcolor}

\geometry{margin=1in}

\hypersetup{
  pdftitle={When a Consciousness Verdict Cannot Be Earned},
  pdfauthor={Daniel Alami},
  pdfsubject={Philosophy of mind, consciousness, identification theory, descent obstruction},
  colorlinks=true,
  linkcolor=blue,
  citecolor=blue,
  urlcolor=blue,
  hypertexnames=false
}

\newtheorem{theorem}{Theorem}
\newtheorem{lemma}{Lemma}
\newtheorem{definition}{Definition}
\newtheorem{corollary}{Corollary}
\newtheorem*{corollary*}{Corollary}
\newtheorem{remark}{Remark}
\newtheorem*{proposition*}{Proposition}

\lstset{
  basicstyle=\footnotesize\ttfamily,
  breaklines=true,
  frame=tb,
  framerule=0.4pt,
  aboveskip=1em,
  belowskip=1em,
  language=Python,
  keywordstyle=\color{blue},
  commentstyle=\color{gray},
  showstringspaces=false
}

\title{When a Consciousness Verdict Cannot Be Earned:\\
\large A Descent-Theoretic Admissibility Criterion}

\author{Daniel Alami}

\date{June 2026}

\begin{document}
\maketitle

%%─────────────────────────────────────────────────────────────────────────────
\begin{abstract}

A claim that a system is or is not conscious is, in the end, a claim about
evidence, only as good as the evidence's power to tell consciousness apart from
its absence. This is an identification problem. A verdict is licensed only when
the target property can be recovered from a declared observation site, and where
it cannot, more data does not help. For a structurally specified class of
hidden-variable systems this recovery fails \emph{necessarily}, surviving every
admissible revision of the measurement language, not only the instrument in
hand. The failure takes a single form, \emph{revision-stable non-effective
descent}: the target cannot be assembled from any admissible family of local
measurements. The underlying factorization is proved for the finite
deterministic and exact finite stochastic cases, with both factorization
regimes, the finite kernel-pair lemma, and conservative invariance mechanically
verified in Lean. The conservative-invariance result shows that no
Morita-equivalent re-presentation of the measurement language removes the
obstruction. The general reduction, that every admissible identification
procedure factors through effective descent, is isolated as an explicit open
obligation, not assumed. The criterion's normative force is confined to
fail-closed audit settings: where it holds, a low-concern verdict is not
licensed, and the default is restriction. Non-identification is not thereby
inadmissible for science in general. Prior category-theoretic work characterizes
the structure of consciousness constructively. This paper runs the other way,
giving a limitative criterion for when an attribution can be made at all. The
criterion is worked through a minimal broadcast-report architecture, where its
global-identity target lands in the $Z = 1$ class while access stays identified,
and through an open-weight transformer. Explicit site constructions for full-scale
Integrated Information, Global Workspace, and Higher-Order theories remain open.
\end{abstract}

\noindent\textbf{Keywords:} consciousness; scientific admissibility;
identification; descent theory; artificial intelligence systems; philosophy of
science

\section*{Technical summary}
A hidden-variable system $H$ generates revision-stable necessary
non-identification of a target property $\Theta$, under the descent-factorization
obligation (Theorem~\ref{thm:descent}, proved here for the finite cases, open in
general), when the target fiber $F \to H$ has non-effective descent over the
admissible observation site $\mathsf{C}_Y$ for every admissible structural
revision $r \in \mathrm{Rev}_\Omega(Y)$.
The revision envelope $\mathrm{Rev}_\Omega(Y)$ covers future structural changes
to the observation framework, not only refinements within a fixed language.
Theorem~\ref{thm:descent} recovers the kernel-pair obstruction lemma
(Lemma~\ref{thm:kernelpair}) as the classical special case.
Integrated Information Theory \citep{tononi2016}, Global Workspace Theory
\citep{baars1988,dehaene2011}, and Higher-Order Theories
\citep{rosenthal2005} are discussed as illustrative sketches of three positions
relative to the criterion, with the explicit site and target fiber for each
left open.

\tableofcontents
\newpage

%%─────────────────────────────────────────────────────────────────────────────
\section{Identification failure as a structural constraint}
\label{sec:intro}

Suppose an AI company claims its model deserves a ``low concern'' verdict on
moral status, meaning there is no significant probability the system is
conscious. What evidence would settle this? The company offers behavioral
evidence (the system does not exhibit distress signals), functional evidence
(it lacks the neural correlates of consciousness), and architectural evidence
(its computation is unlike human cognition). A regulator reads the evidence and
asks whether it separates consciousness from an observational duplicate.

None of this connects the target to the evidence in the right way. A system
that lacks consciousness can still duplicate behavioral competence.
Neural-correlate evidence is tied to human neural architecture. Architectural
dissimilarity shows that two systems differ, while leaving open whether either
of them lacks consciousness. So the proposed tests leave two stories standing,
and no experiment on the table separates them.

Philosophers call this the ``hard problem.'' Stated more precisely, it is an
\emph{identification problem}. A parameter is identified when different true
values yield different observable distributions. Once identification fails, more
data does not help.

The central question is whether identification failure for consciousness holds
\emph{necessarily}, surviving admissible improvements in measurement, or only
\emph{contingently}, so that a better admissible test could remove it. The
answer decides whether a failed test should prompt further measurement or a
restricted verdict.

For a well-defined class of systems the failure is necessary
relative to the declared observation site and revision envelope, conditional on
the descent-factorization obligation stated in Theorem~\ref{thm:descent} and
proved below for the finite deterministic and exact finite stochastic cases.
A system in this class cannot receive an auditable low-concern verdict without
a successful identification certificate.

%%─────────────────────────────────────────────────────────────────────────────
\section{Formal preliminaries}
\label{sec:prelim}

``Admissible observations'' names
the measurements permitted for the system under evaluation. Two hidden states
are experimentally indistinguishable when they agree under all of those
observations. And the \emph{experimentability witness} records the minimal
typed specification a scientific consciousness claim must provide. A theory
that cannot supply this witness has made no falsifiable claim, so the theorems
below leave it untouched.

For the setup, let $H$ be the set of possible hidden states of a system (its
internal state
space: neurons firing, weights activating, whatever physical process underlies
behavior).
Let $\Theta: H \to T$ be the target map (e.g., ``is this state conscious or
not?'').
Let $Y$ be the collection of admissible observations: every measurement,
perturbation, imaging technique, and behavioral test permitted.
Let $E_Y: H \to \mathcal{R}_Y$ be the observation function mapping hidden
states to observable records.

\begin{definition}[When two hidden states look identical]
\label{def:congruence}
Hidden states $h_1$ and $h_2$ are \emph{experimentally indistinguishable} if
every admissible measurement gives the same distribution:
\[
h_1 \equiv_Y h_2
\;\iff\;
\sup_{e \in \mathrm{Eff}(Y)} |p(e \mid h_1) - p(e \mid h_2)| = 0.
\]
On classical deterministic systems this simplifies to
$E_Y(h_1) = E_Y(h_2)$.
\end{definition}

\begin{definition}[Experimentability witness]
\label{def:witness}
A theory $(H, Y, E_Y, \Theta)$ is \emph{experimentable} if it supplies a
typed specification
$W = (\mathcal{C}, \mathcal{R}, \mathcal{T}, H, Y, E_Y, \Theta)$
in which $\mathcal{R}$ and $\mathcal{T}$ have decidable equality.
By ``decidable equality'' we mean a procedure that definitively answers
``did these two records match?'' and ``did these two target values match?''
Absent such a specification, the theory makes no falsifiable scientific claim
in the sense used here.
\end{definition}

\begin{remark}
Faced with a theory that cannot supply the witness, a governance regulator
should require the witness before deliberating about the theory at all. IIT,
GWT, and HOT occupy different positions under this test, as
Section~\ref{sec:implications} shows: IIT does not yet supply the witness for
full $\Phi$-based identification, whereas GWT and HOT can supply witnesses for
narrower targets.
\end{remark}

%%─────────────────────────────────────────────────────────────────────────────
\section{Structural characterization of the non-identification class}
\label{sec:characterization}

The hidden-variable systems that fall in the necessary non-identification class
are picked out by four structural properties, none of them mentioning
consciousness. A system satisfying all four generates $Z = 1$ by the theorem in
\S\ref{sec:descent}, whose machinery \S\S\ref{sec:kernelpair}--\ref{sec:descent}
develop.

\begin{definition}[P1, local contextual presentability]
A family of admissible local contexts $\{U_\alpha \to H\}$ exists such that
each $U_\alpha$ supports local target assignments.
The target $\Theta$ is partially evaluable at the level of individual contexts,
even when not globally evaluable.
\end{definition}

\begin{definition}[P2, nontrivial transition structure]
On overlapping contexts $U_\alpha \cap U_\beta$, local target assignments are
related by transition maps
$g_{\alpha\beta} : F_\beta|_{U_{\alpha\beta}} \to F_\alpha|_{U_{\alpha\beta}}$
that satisfy local compatibility on every pair but fail global compatibility
across all contexts simultaneously.
\end{definition}

\begin{definition}[P3, nontrivial holonomy]
The context groupoid $\Pi_Y(H)$ acts on the target fiber $F$ with no admissible
global fixed section: $\mathrm{Fix}(\rho) = \emptyset$.
The descent datum is non-effective: no global section is compatible with all
local data simultaneously.
\end{definition}

\begin{definition}[P4, equivariant admissible operations]
Every admissible intervention or observation refinement $Y' \in \mathrm{Adm}(Y)$
preserves the holonomy structure:
$\rho_{Y'} \circ a = a \circ \rho_Y$ for every admissible action $a$.
No expansion or refinement of the observation family dissolves the obstruction.
\end{definition}

Properties P1--P3 fix a static obstruction at one observation level. P4 makes
the obstruction persist, since it guarantees that the holonomy survives every
admissible intervention. Where P4 fails, the obstruction is contingent, and a
better instrument can in principle remove it. Where all four hold, the
obstruction is necessary, beyond the reach of any improvement in measurement
methodology.

These properties depend only on the topology of the causal graph and on how the
target fiber descends over the observation site. Substrate, material
composition, and biological origin drop out. Any physical system, biological or
digital, that instantiates equivalent P1--P4 structure falls in the same
identification class.

P1--P4 need not exhaust every possible source of the hard problem. For a system
with P1--P4, conservative expansions of the
structural language draw no new structural distinctions, so they cannot supply
the sufficiency bridge from non-identification to consciousness attribution, and
the obstruction is stable under conservative re-presentation.
Non-conservative expansions can matter, yet they change the admissible site $Y$
and must be declared in the revision envelope $\Omega$. The boundary result
below concerns this conservative-stability claim, and a complete taxonomy of
possible consciousness obstacles lies outside it.

%%─────────────────────────────────────────────────────────────────────────────
\section{The kernel-pair non-factorization lemma}
\label{sec:kernelpair}

Identification can fail in a way that is at once the simplest and the easiest to
miss. Two hidden states look identical from outside, since every measurement
returns the same result, yet they differ in whether $\Theta$ holds. In the
consciousness case, a conscious system and a non-conscious system produce
identical observable records under every admissible test, so no experiment
distinguishes them. This is the kernel-pair obstruction.

\begin{definition}[Kernel pair]
\label{def:kernelpair}
The \emph{kernel pair} of $E_Y$ is the set of all experimentally
indistinguishable pairs:
\[
K_Y = \{(h_1, h_2) \in H^2 : h_1 \equiv_Y h_2\}.
\]
Think of $K_Y$ as the observation window's blind spot.
The quotient $H/\!\sim_Y$ is the compressed state space in which
indistinguishable states are treated as one.
\end{definition}

\begin{lemma}[Kernel-pair obstruction]
\label{thm:kernelpair}
For an experimentable theory $(H, Y, E_Y, \Theta)$:

\medskip
\noindent\textit{$\Theta$ is necessarily non-identified if and only if there
exist hidden states $h_1$, $h_2$ that look identical from outside
($h_1 \equiv_Y h_2$) but differ in the target
($\Theta(h_1) \neq \Theta(h_2)$).}

\medskip
\noindent Equivalently: $\Theta$ does not factor through the experimentally
compressed state space $H/\!\sim_Y$.
\end{lemma}

\begin{proof}[Proof sketch]
If the pair exists, no estimator on observable records can separate $h_1$
from $h_2$.
If no such pair exists, $\Theta$ is constant on each indistinguishability
class and is identified.
\end{proof}

This is the standard non-identification condition of \citet{pearl2009},
instantiated in the present typed setup.

A concrete check follows. Build a six-state system
$\{a, b, c, d, e, f\}$.
Let the observation function collapse $\{a, b\}$ to the same record.
Set $\Theta(a) = 0$, $\Theta(b) = 1$.
The kernel pair $(a,b)$ fires: identical observation, different target, with
residual ambiguity $\log_2 2 = 1.000$ bit, and no symmetry or automorphism
enters. The argument turns on plain indistinguishability, with nothing
gauge-theoretic about it.

The lemma falls short in one respect. Lemma~\ref{thm:kernelpair} needs the
compressed state space $H/\!\sim_Y$ to be a well-behaved mathematical object. On
quantum systems, systems with non-abelian symmetries, or infinite-dimensional
spaces, the quotient may fail to be constructible. Theorem~\ref{thm:descent}
handles all of these.

%%─────────────────────────────────────────────────────────────────────────────
\section{The descent-obstruction theorem}
\label{sec:descent}

Lemma~\ref{thm:kernelpair} catches the case where two states look identical and
the target differs. A more general obstruction lurks beyond it. Every
\emph{pair} of measurement contexts can be locally consistent while no single
global assignment fits all contexts at once. Quantum mechanics knows this
phenomenon as contextuality, and it carries over to consciousness attribution.
Theorem~\ref{thm:descent} states the local-to-global obstruction at the level of
effective descent.

\subsection{Limits of \v{C}ech cohomology}
\label{sec:cech}

\citet{abramsky2011} developed the standard tool for contextuality, sheaf
cohomology, for quantum systems. A nonzero cohomology class
$\check{H}^1(\mathcal{U}_Y, \hat\Theta) \neq 0$ certifies that local target
assignments cannot be glued globally, detecting cases that
Lemma~\ref{thm:kernelpair} misses.

Yet $\check{H}^1$ does not tell the whole story. It demands a good cover
(geometric regularity), abelian target structure (with addition and subtraction
defined), and an obstruction lodged in exactly degree~1. Quantum POVM systems,
non-abelian symmetry groups, and systems with non-standard topologies all
violate one or another of these conditions.

The kernel-pair case (Theorem~\ref{thm:kernelpair}) is a direct counterexample,
with $\check{H}^1 = 0$ but $D = 1$. A detector that checks only $\check{H}^1$
recovers half the cases on a minimal two-case suite, recall $= 1/2$, which is
too coarse to rely on.

\subsection{Effective descent}
\label{sec:descent-intuition}

Imagine reconstructing the value of $\Theta$ by combining local measurements
from overlapping contexts. \emph{Effective descent} means that some procedure
takes all the local results, checks them for overlap consistency, and assembles
them into a single global answer. When descent fails to be effective, no such
assembly exists: every pair of local measurements can agree, and the global
assembly still goes through nowhere.

\begin{definition}[Descent-obstruction flag]
For admissible observation family $Y'$:
\[
D_{Y'}(X) =
\begin{cases}
1 & \text{if } F|_{\mathsf{C}_{Y'}} \text{ has non-effective descent} \\
0 & \text{otherwise.}
\end{cases}
\]
\end{definition}

\begin{definition}[Revision envelope]
\label{def:revenv}
Let $\Omega$ denote the declared physical, computational, and operational
constraints on what structural changes to the observation/intervention
framework are admissible.
A \emph{structural revision} $r$ is any change to the admissible site that
falls within $\Omega$.
$\mathrm{Rev}_\Omega(Y)$ denotes the set of such revisions.
A revision is \emph{conservative} (Morita-equivalent) if it induces an
equivalence of structural model categories over the same admissible site,
preserving the target fiber, covering families, and context structure.
A revision is \emph{non-conservative} if it adds a genuinely new admissible
observable, intervention, or target-relevant operation not definable from
the old site.
\end{definition}

\begin{definition}[Revision-stable persistent non-effective descent]
\label{def:revstable}
The identification obstruction is \emph{revision-stable} if it survives every
admissible structural revision:
\[
Z = f(X, Y) = \min_{r \in \operatorname{Rev}_\Omega(Y)} D(r_* X) = 1.
\]
When $Z = 1$, no measurement protocol, refinement, expansion, or structural
revision of the observation family within $\Omega$ lets the target be assembled
globally from local data. The earlier formula
$\min_{Y' \in \mathrm{Adm}(Y)} D_{Y'}(X)$ is the special case in which $\Omega$
holds only refinements within a fixed structural language.
Definition~\ref{def:revstable} reaches all admissible structural revisions.
\end{definition}

\begin{proposition*}[Conservative Invariance]
If $e: \mathcal{L}_D \to \mathcal{L}_D^+$ is a Morita-conservative expansion
over the admissible site (the induced functor on structural model categories
is an equivalence compatible with the site topology and target fiber), then:
\[
F \text{ has effective descent in } \mathcal{L}_D
\iff
e_* F \text{ has effective descent in } \mathcal{L}_D^+.
\]
\end{proposition*}

\begin{proof}[Proof sketch]
Conservative presentations induce equivalent descent categories only when the
equivalence preserves the site topology: covering families, pullback
compatibility, and the target fiber are transported together.
Equivalences preserve existence and uniqueness of global objects.
Effective descent is exactly existence of a global object realizing the local
datum.
No identification-failure premise is used.
\end{proof}

\begin{corollary*}
Conservative expansion cannot overturn $Z = 1$.
If the target fiber has non-effective descent, no Morita-equivalent
re-presentation of the structural language changes this.
Any future argument that a richer structural language supplies the sufficiency
bridge must add a non-conservative admissible operation. Such an operation
changes $Y$, and $\Omega$ must declare it. Treating it as a language-only
clarification will not do.
\end{corollary*}

\subsection{The descent criterion}
\label{sec:maintheorem}

\begin{proposition*}[Finite deterministic factorization]
Let $H$ and $\mathcal{R}_Y$ be finite sets, let $E_Y:H\to\mathcal{R}_Y$ be a
deterministic observation map, and let $\Theta:H\to T$ be a target with
decidable equality.
An identification procedure using only admissible records exists if and only
if there is a map $\widehat{\Theta}:\mathcal{R}_Y\to T$ such that
\[
\Theta=\widehat{\Theta}\circ E_Y.
\]
Equivalently, $\Theta$ is constant on every fiber of $E_Y$.
\end{proposition*}

\begin{proof}
If such $\widehat{\Theta}$ exists, the procedure observes $r=E_Y(h)$ and
returns $\widehat{\Theta}(r)$.
If an identification procedure using only records exists, its output is a
function of $r$ alone.
Two hidden states in the same fiber of $E_Y$ therefore receive the same output,
so correctness requires $\Theta$ to be constant on that fiber.
Define $\widehat{\Theta}(r)$ to be that common value for each realized record.
Unrealized records can be assigned arbitrarily.
This is the finite deterministic form of effective descent: local record data
assemble a unique global target value exactly when the target factors through
the record object.
\end{proof}

\begin{proposition*}[Exact finite stochastic factorization]
Let $H$ and $\mathcal{R}_Y$ be finite sets.
Let $K:H\to\Delta(\mathcal{R}_Y)$ assign each hidden state an exact observation
distribution over records, and let $\Theta:H\to T$ be a target with decidable
equality.
An exact identification procedure whose information is exhausted by the
stochastic observation site exists if and only if $\Theta$ factors through the
kernel map:
\[
\Theta=\widehat{\Theta}\circ K
\]
for some map $\widehat{\Theta}$ defined on the image of $K$.
Equivalently, whenever $K(h)=K(h')$, $\Theta(h)=\Theta(h')$.
\end{proposition*}

\begin{proof}
If $\Theta$ factors through $K$, the procedure reads the exact kernel $K(h)$
and returns $\widehat{\Theta}(K(h))$.
Conversely, any exact procedure using only the stochastic site must assign the
same answer to hidden states with the same observation distribution.
Correctness therefore requires $\Theta$ to be constant on each fiber of $K$.
Define $\widehat{\Theta}$ on each realized kernel by that common value.
\end{proof}

This proposition concerns exact distributional information.
It does not say that one finite sample identifies the target.
If two kernels overlap, one-sample classification can have nonzero Bayes error
even when the exact kernels are distinct.
A sharper obstruction is at work here: if the kernels are identical and the
targets differ, no amount of sampling from that site identifies the target.

\begin{sloppypar}
Grothendieck's theory \citep{grothendieck1971sga} establishes that
effective descent over a site is exactly the condition under which coherent
local data determine a global object. The propositions above prove the reduction
for finite deterministic systems and exact finite stochastic systems.
Every identification procedure we know of (maximum-likelihood estimation,
instrumental variables, Pearl's do-calculus \citep{pearl2009}, Bayesian
updating) takes the same form: given local measurement records that cover the
state space at some observational resolution, recover the target parameter. That
recovery has the structure of effective descent of the target fiber over the
measurement site. So when the target fiber has non-effective descent, no
reconstruction procedure of this form should succeed, because the local data
carry too little information to assemble a global answer. The
reduction is argued here and proved for the formal components named in
\S\ref{sec:limits}. A full proof that every admissible identification procedure
in every substrate class factors through effective descent lies beyond this
article.
\end{sloppypar}

Fix a site $(\mathsf{C}_Y, J)$, where $J$ is a Grothendieck topology on the
category of admissible measurement contexts. Non-effective descent of the target
fiber $F$ over this site says that no compatible family of local sections
uniquely assembles into a global section. One invariant, non-effective descent
of the target fiber, runs across every substrate, even though the
computational tool that detects it shifts with the topology. Classical finite
systems call for \v{C}ech $H^1$, quantum POVM systems for the
Abramsky--Brandenburger empirical-model method \citep{abramsky2011}, and
coalgebraic systems for bisimulation-invariant descent. A substrate-specific
table below lists these implementations. The claim that each substrate
instantiates the same descent criterion under its natural topology is a scoped
mathematical program, one this article does not complete.

\begin{theorem}[General descent-obstruction criterion, conditional on the descent-factorization obligation]
\label{thm:descent}
Let $(H, Y, E_Y, \Theta)$ be an experimentable theory with declared revision
envelope $\Omega$. Suppose the descent-factorization obligation holds for the
substrate: admissible identification procedures factor through effective descent
of the target fiber over the observation site. This obligation is proved for the
finite deterministic and exact finite stochastic cases above. It remains open in
general (\S\ref{sec:limits}) and is assumed established nowhere else in this
paper.

\medskip
\noindent\textit{Then $\Theta$ is necessarily non-identified from admissible
observations if and only if $Z = 1$.}
\end{theorem}

\begin{proof}[Proof sketch under the descent-factorization obligation]
If $Z = 1$, then every admissible revision $r \in \mathrm{Rev}_\Omega(Y)$
gives non-effective descent.
No reconstruction of $\Theta$ from observable records exists for any
admissible protocol or structural revision within $\Omega$, since descent is
the most general reconstruction principle and the Proposition (Conservative
Invariance) rules out conservative expansion as a route around the
obstruction.
If $Z = 0$, some revision $r$ gives effective descent, so the target fiber
can be reconstructed from local data via the descent morphism.
\end{proof}

The non-applicability boundary can be stated formally.
A substrate is in-domain (the theorem applies) if and only if all four of the
following are defined: (1)~an admissible context site $\mathsf{C}_Y$,
(2)~a target-fiber object $F \to H$, (3)~descent semantics for that substrate
category, and (4)~a declared revision envelope $\mathrm{Rev}_\Omega(Y)$.
If any term is absent, the correct output is non-applicability (neither
$Z = 0$ nor $Z = 1$) and the governance default is fail-closed.

\subsection{The sufficiency boundary}
\label{sec:sufficiency}

Theorem~\ref{thm:descent} says when P1--P4 systems lie in the non-identification
class under the descent-factorization obligation. Whether every P1--P4 system is
conscious, the sufficiency direction, we settle here as a formal boundary
result.

Conservative expansions supply no sufficiency. Suppose a future structural
language $\mathcal{L}_D^+$ is conservative over $\mathcal{L}_D$
(Morita-equivalent, by the Proposition above), so that it draws no new structural
distinctions between models:
\[
\mathrm{Mod}(\mathcal{L}_D^+) \simeq \mathrm{Mod}(\mathcal{L}_D)
\implies \Delta(\text{distinguishable structural cases}) = 0.
\]
So a conservative expansion can turn
$\mathrm{P1\text{--}P4} \not\Rightarrow \Theta$ into
$\mathrm{P1\text{--}P4} \Rightarrow \Theta$ only if the implication was already
derivable from the old structure, which the current framework nowhere
licenses. Conservative re-presentations of the language leave the sufficiency
bridge unbuilt.

Non-conservative expansions change the domain. Suppose
$\mathcal{L}_D^+$ adds a primitive that directly tracks $\Theta$. That clarifies
no language, since it adds a new admissible probe and changes
$Y$. A sufficiency bridge derived from $\mathcal{L}_D^+$ would then follow from
the new
$Y$, with the old structural framework doing none of the work. Any such
expansion must be declared in $\Omega$ and evaluated as a non-conservative
revision.

So the negative sufficiency result, that P1--P4 systems are not thereby shown to
be conscious, owes nothing to an impoverished structural language. It holds
under every conservative expansion, and only a declared non-conservative
revision that changes what is admissible can defeat it. Within any
observable-relative framework, the bridge from the non-identification class to
consciousness attribution calls for an exogenous commitment that lies outside
the framework. Other axiomatizations may add further commitments. The scope here
is only what observable-relative formal frameworks can determine before any
such commitment is made.

\subsection{Comparison of approaches}

Table~\ref{tab:coverage} summarizes how the kernel-pair lemma, cohomological
diagnostics, and the descent theorem cover the representative substrate
classes used in the paper.

\begin{table}[H]
\centering
\renewcommand{\arraystretch}{1.3}
\begin{tabular}{lcccc}
\toprule
Case & Lemma~\ref{thm:kernelpair} & $\check H^1$ & Thm.~\ref{thm:descent} \\
\midrule
Two states look identical, $\Theta$ differs & catches & misses & catches \\
Contextuality triangle (classical)           & misses  & catches & catches \\
Quantum POVM contextuality                   & n/a     & catches & catches \\
Non-abelian symmetry system                  & partial & partial & catches \\
Identifiable control                         & silent  & silent  & silent  \\
\bottomrule
\end{tabular}
\caption{Coverage across substrate classes.
Lemma~\ref{thm:kernelpair} is the $H^0$ special case of
Theorem~\ref{thm:descent}.
$\check{H}^1$ is a sufficient diagnostic under classical good-cover
conditions only.}
\label{tab:coverage}
\end{table}

In practice the right mathematical tool depends on the substrate:

\begin{table}[H]
\centering
\renewcommand{\arraystretch}{1.3}
\begin{tabular}{lp{8cm}}
\toprule
Substrate & Use this obstruction tool \\
\midrule
Classical finite systems      & \v{C}ech $H^1$; CSP global-section check \\
Quantum (POVM) systems        & Abramsky--Brandenburger empirical model \citep{abramsky2011} \\
Non-abelian systems           & Non-abelian $H^1$, torsors, groupoid descent \\
Measurable/probabilistic      & Invariant $\sigma$-algebras, measurable descent \\
Coalgebraic / process systems & Bisimulation-invariant descent \\
\bottomrule
\end{tabular}
\caption{Substrate-specific tools for computing the descent-obstruction flag.}
\label{tab:substrates}
\end{table}

Appendix~\ref{app:suite} and the experiments directory make the checks below
executable.
\begin{itemize}
\item \textit{Contextuality triangle}: $\delta^0$ rank over $\mathbb{F}_2 = 2$;
  $\dim\check{H}^1 = 1$; zero global sections; $D = 1$.
\item \textit{Kernel-pair case}: $\check{H}^1 = 0$; $D = 1$;
  $H^1$-only recall on two-case suite $= 1/2$.
\item \textit{Identifiable control}: $D = 0$.
\item \textit{Finite stochastic kernel pair}: identical observation kernels
with different targets fail exact factorization; distinct kernels pass exact
factorization while retaining finite-sample error.
\item \textit{Boolean recurrent architecture}: report-only traces leave a
latent target unidentified across all checked input sequences; a direct latent
reader removes the obstruction by changing the site.
\end{itemize}

In the consciousness applications below, the simpler branches do most of the
work. IIT fails already at the witness stage, while GWT and higher-order
theories raise kernel-pair worries whenever access or metacognitive evidence
serves a phenomenal target. The descent generalization
comes into its own on contextual, quantum-adjacent, and non-abelian substrates,
where those simpler diagnostics no longer cover the space of identification
failures.

%%─────────────────────────────────────────────────────────────────────────────
\section{The governance corollary}
\label{sec:governance}

When $Z = 1$ for a system, no experiment (no certificate, audit, or disclosure
protocol) can establish that $\Theta$ is identified. An epistemic claim and a
normative premise follow, and they stay apart.

The epistemic claim flows from Theorem~\ref{thm:descent} alone. A
governance protocol that issues a low-concern verdict on consciousness for such
a system issues a verdict it cannot substantiate, because the theorem shows that
no admissible experiment can supply the required identification evidence.

The normative premise: a low-concern verdict requires
epistemic support, and where no support is possible the appropriate governance
default is to restrict. The premise earns its place from the asymmetry of error
costs under consciousness uncertainty, which warrants precautionary closure.
Being conditional, it dominates only when a mistaken low-concern verdict would be
severe, hard to reverse, and borne by a possible subject, while the cost of
withholding the verdict stays bounded by delay, further disclosure, narrower
deployment, or a graduated precautionary regime. Where those conditions fail,
the mathematics still delivers the epistemic negative (the low-concern verdict
lacks identification evidence) while leaving the practical policy response open.
Given the premise, restriction follows as a theorem. Without it, the mathematics
still shows that low-concern verdicts are epistemically unwarranted, yet says
nothing about the governance response. The criterion's normative force stays
confined to fail-closed audit and governance settings. It makes no claim that
non-identification is inadmissible for science in general.

\begin{corollary}[Fail-closed verdict rule]
\label{cor:governance}
If $Z = f(X,Y) = 1$, then for every admissible $Y' \in \mathrm{Adm}(Y)$:
\[
\Theta \notin \mathrm{Desc}(Y')
\quad\implies\quad
\text{low-concern verdict on } \Theta \text{ is not admissible.}
\]
\end{corollary}

\begin{proof}
By $Z = 1$ every $Y'$ gives non-effective descent.
A low-concern verdict requires establishing identifiability under some $Y'$.
Non-effective descent entails no such demonstration exists.
\end{proof}

A verdict can come out three ways.
\begin{center}
\renewcommand{\arraystretch}{1.3}
\begin{tabular}{ll}
\toprule
Situation & Correct verdict \\
\midrule
Experimentability witness fails & Non-evaluable $\to$ no low-concern verdict \\
Witness holds and $Z = 1$       & Identification impossible $\to$ verdict formally forbidden \\
Witness holds and $Z = 0$       & Possible $\to$ requires adversarial certificate + replication \\
\bottomrule
\end{tabular}
\end{center}

Both the ``witness fails'' and ``$Z = 1$'' branches are fail-closed, so a
governance protocol open to low-concern verdicts in either branch contradicts
Corollary~\ref{cor:governance}.

Most governance frameworks for AI consciousness are precautionary, advising
caution because the evidence is incomplete. The corollary isolates a narrower
case. When $Z = 1$ relative to a declared $\Omega$, the declared evidence base
simply cannot identify the target. On the normative premise adopted above, a
low-concern verdict then lacks the support it requires.

In its audit form the corollary is direct. A low-concern verdict must state the
site, target, revision envelope, classifier, challenge record, revocation rule,
and provenance record it relies on. Should any required field go missing,
the verdict is not evaluable under this criterion.
As a downstream policy matter beyond this paper, the split between the
access/report target and the phenomenal target supports a deployment license
tied to the access/report verdict alone and barred from carrying that verdict
over to phenomenal status.

\subsection{The disclosure requirement}
\label{sec:disclosure}

A protocol can apply the governance rule only once the system's descent
structure is on the table. A low-concern claim that omits the following six
elements is not evaluable. It receives the same fail-closed verdict as
$Z = 1$.

Any organization asserting a low-concern verdict for $\Theta$ must publish:

\begin{enumerate}
\item The context category/site $\mathsf{C}_Y$: a complete specification of
  admissible measurements and interventions.
\item The target fiber object $F \to H$: a structural specification of
  $\Theta$ as a fiber over the hidden-state space.
\item The obstruction datum $X$: the descent data (transition maps, holonomy,
  kernel-pair constraints, contextual compatibility) that enables a
  third-party auditor to evaluate whether the system falls in the P1--P4
  class.
\item The admissible revision envelope $\Omega$: the declared physical,
  computational, and operational constraints on structural revision, with
  explicit classification of each proposed language expansion as conservative
  (Morita-equivalent) or non-conservative (genuinely new admissible probe or
  operation).
\item A revision-stable descent test result: a demonstration that
  $Z = \min_{r \in \operatorname{Rev}_\Omega(Y)} D(r_* X) = 0$, with a
  certificate replicable by a third party under the declared $\Omega$.
\item A non-applicability statement or confirmation: an explicit declaration
  that all four in-domain conditions (\S\ref{sec:maintheorem}) are satisfied,
  or identification of which condition fails and why non-applicability, not
  $Z = 0$, is the correct output.
\end{enumerate}

Any missing element forces fail-closed. Without $\Omega$, $Z$ is not evaluable
under the revision-stable criterion. Undeclared non-conservative expansions then
let a claimant revise the admissible site after the fact to dodge a $Z = 1$
verdict.

Each element guards against a specific governance failure mode. Undeclared site
data lets a claimant gerrymander the admissible evidence base. Undeclared
holonomy data hides whether the system falls in the P1--P4 class. A missing
$\Omega$ opens the door to strategic omission of the non-conservative revisions
that would resolve the obstruction. And a missing revision-stable certificate
licenses unilateral identifiability claims that no one can audit as the
structure changes.

\subsection{\texorpdfstring{The $\Omega$-audit protocol}{The Omega-audit protocol}}
\label{sec:omega-audit}

Naming $\Omega$ as an input, as the disclosure requirement does, leaves it open
to gaming. A claimant who wants a low-concern verdict can shape the envelope
before the descent test runs, and that choice can decide whether $Z$ is
evaluable at all. An $\Omega$ declaration therefore needs its own audit
protocol.

The protocol is organized around six failure modes. They need not exhaust the
space. They form a finite audit framework against which a later completeness
theorem could be tested.

\emph{Under-specification.} A claimant may declare $\Omega$ vaguely, for example
by saying that admissible revisions are those consistent with current physics or
current practice, which leaves later challengers to be accepted or excluded
after the fact. To remedy this, enumerate finitely. The declaration must list
admissible revision classes (new observable types, intervention types,
model-access regimes, target-relevant operations, and structural vocabulary
changes), each carrying an inclusion or exclusion rationale. Anything outside the
enumeration is inadmissible for the current certificate, though the enumeration
itself stays open to public challenge.

\emph{Strategic exclusion.} A claimant may exclude a revision class that would
resolve the obstruction by calling it physically impossible, operationally
infeasible, or out of scope. An exclusion is legitimate only when the reason is
typed and auditable. A finite list of reason types is admissible: physical
impossibility, computational intractability, regulatory prohibition, or cost
above a declared threshold, where the threshold itself is part of the claim. Let
an excluded class lack one of these reasons, or let the reason fail under audit,
and the low-concern verdict is revoked or reduced to non-applicability.
Each exclusion rationale is negative, recording why a class stays out. The
claimant must also give its positive counterpart, the minimal set of
non-conservative probes that would defeat the obstruction and return $Z = 0$ if
they were admissible and run. Fixing that set before certification blocks a later
denial that a surfaced probe was ever relevant. This does not certify
completeness, since the falsifying set need not be finitely nameable in advance,
but it closes one route to revising the claim after a challenge.

\emph{Misclassification.} A claimant may add a probe that changes the admissible
site while calling it a conservative language expansion. The audit rule is
simple: any expansion claimed as conservative must supply a witness that the new
presentation preserves the target fiber, covering families, and context
structure. In mathematical cases this is a Morita-equivalence witness, and in
applied cases the corresponding operational equivalence argument. With no
witness, the expansion counts as non-conservative and triggers a fresh $Z$
evaluation on the changed site.
For finite structural fragments, where the equivalence is decidable, the witness
must be machine-checkable, as the finite-alias case already is in Lean
(\S\ref{sec:limits}), so a third party can recompute the equivalence and need not
take the claimant's reading on trust. The equivalence is undecidable in general,
and there the witness reverts to a refereed proof, with the fail-closed default
holding when none is supplied.
That referee pool is part of the apparatus the provenance certificate covers.
Fail-closed also blunts an attack by attrition: a claimant who floods the review
with proofs too large to check leaves the referees unable to accept the witness,
so the default holds and the expansion counts as non-conservative. Only an
affirmative, captured certification can pass a site-changing probe off as
conservative. Referee exhaustion cannot.

\emph{Post-hoc revision.} A claimant may alter $\Omega$ after seeing data or
after receiving a challenge. Here the remedy is preregistration discipline. The
envelope must be declared before certification, timestamped, and hashed. Any
later change starts a new certificate. The earlier verdict inherits no evidence
from the revised site until the new evaluation runs.

\emph{Residual discovery failure.}
A claimant may satisfy the listed declaration rules while a material
non-conservative probe class is never surfaced. Seeded omission tests do not
solve this. They can show that a red team detects known planted defects under
benchmark conditions, but they cannot estimate how likely an adversarially
chosen future probe class is to remain outside the audit community's search
space. So seeded omissions function as a skill gate, leaving residual closure
uncertified. For operational use, write $\varepsilon_{\mathrm{esc}}$ for the
residual adversarial escape load and $\iota_{\mathrm{audit}}$ for verified audit
integrity: independent selection, escrow funding, access parity, public
challenge logs, conflict checks, and whistleblower channels.
The admissibility score has the form
\[
\sigma_{\mathrm{adm}} = \iota_{\mathrm{audit}}(1-\varepsilon_{\mathrm{esc}}).
\]
Whenever a required dossier field is absent, access parity fails, red-team
independence is compromised, or $\varepsilon_{\mathrm{esc}}$ cannot be bounded
under the declared adversary model, the output is non-applicability. A
seeded-omission success rate is no admissible bound for
$\varepsilon_{\mathrm{esc}}$.

One concrete instantiation of the declared adversary model is an escrow stake. A
claimant seeking a low-concern verdict places capital or compute in escrow that
funds independent red teams, who earn the stake by exhibiting an admissible probe
that flips the verdict within the declared rules, which then revokes the
certificate. This sharpens discovery by giving the search real resources and a
real incentive. It does not certify closure. A quiet escrow period, with no
probe found, shows only that none was found under that search and is no bound on
$\varepsilon_{\mathrm{esc}}$, so the residual-closure requirement above still
holds.

\emph{Apparatus provenance failure.} A claimant may shape the people or
infrastructure that generate the challenge record before the public record is
fixed. The visible dossier can still display formal independence, escrow
funding, access parity, public logs, and a timestamped envelope, even as the
hidden dependence sits upstream of all those records. The channels at issue
include funding, employment prospects, advisory ties, compute or API access,
benchmark infrastructure, data access, publication access, shared investors, and
informal reciprocal arrangements. Write this residual as $X_{\mathrm{prov}}$. A
low-concern verdict therefore calls for a provenance certificate covering
auditors, red teams, classifiers, panel members, infrastructure providers, and
benchmark or data providers. The certificate must aggregate claimant-originated
benefits by source class. Per-edge materiality thresholds may cap how much
detail is published, yet they shelter no one from aggregation. Where
source-class provenance cannot be certified, the provenance gate fails and the
output is non-applicability.

All of this shifts the burden of a low-concern verdict onto the claimant, who
must publish the envelope before the descent test, defend each excluded revision
class, supply equivalence witnesses for conservative expansions, certify the
provenance of the challenge apparatus, preserve challenger independence, and
accept revocation when the envelope was too vague or moved after inspection. A
third-party auditor need not settle consciousness. The auditor asks only whether
a fixed and contestable revision envelope supports the claimed low-concern
verdict.

Those six failure modes are the main attacks this proposed protocol blocks.
A later reduction theorem would show that every strategy for gaming an
$\Omega$ declaration either leaves the envelope under-specified, strategically
excludes a relevant revision class, misclassifies a site-changing operation as
conservative, revises the envelope after evidence is observed, exploits a
failure of the discovery apparatus, or contaminates the provenance of the
apparatus that generates the public record.
No such reduction theorem is claimed here.
If such a reduction holds, the audit protocol is complete relative to the
declared site and target fiber.
If it fails, the counterexample identifies another failure mode and the
protocol must be expanded.
Hardest of all is the candidate counterexample of a future probe class that
cannot be named at declaration time.
A completeness proof must show that this reduces to one of the public-record
failure modes, or else treat it as residual discovery failure.
Operationally, this means that a low-concern dossier must not infer residual
closure from planted omissions.
It must either bound residual escape under a declared adversary model and a
public challenge process, or fail closed.
A second counterexample is provenance atomization, where a claimant splits
influence across sub-threshold proxy paths so that no single edge looks material
even as the aggregate claimant-originated exposure is material. The protocol
blocks this only when the provenance certificate aggregates by source class. A
five-mode reduction that files all such cases under residual discovery sweeps too
widely, classifying by failed outcome where it should classify by a
pre-specified mechanism.
A reduction theorem over these six public objects is the natural formal target
for follow-on work: declaration content, exclusion boundary, conservative
classifier, certification timing, discovery apparatus, and apparatus
provenance.

%%─────────────────────────────────────────────────────────────────────────────
\section{A worked application: a minimal broadcast-report architecture}
\label{sec:worked}

Consider a small cognitive architecture. It is not a model of
biological consciousness. It shows how to specify the site, target fiber,
holonomy, and P4 condition in a case where one can check the obstruction by
inspection.

The architecture has a perceptual store $P$, a workspace $W$, and a report
controller $R$.
The hidden state $h \in H$ contains a latent episode token and three local
encodings of that token, one in each module.
The admissible observation site $\mathsf{C}_Y$ has three contexts:
\[
U_{PW},\qquad U_{WR},\qquad U_{RP}.
\]
Each context exposes pairwise agreement records between two modules.
No admissible probe reads all three encodings at once.
This restriction is an ordinary instrumentation limit. The available protocols
can test pairwise access, workspace availability, and report alignment, yet no
single record identifies the latent episode token across all modules at once.

The target $\Theta$ is global episode identity: whether the three local
encodings are assignments of one and the same latent episode.
The target fiber $F \to H$ carries a local binary label for the episode token
in each context.
On overlaps, local labels are related by transition maps.
Take two transition maps to preserve the label and one to flip it:
\[
g_{PW,WR} = +1,\qquad
g_{WR,RP} = +1,\qquad
g_{RP,PW} = -1.
\]
The product around the cycle is $-1$. Every pairwise comparison is locally
satisfiable, yet no consistent global label survives across all three contexts.
The target fiber therefore has non-effective descent over $\mathsf{C}_Y$, and
Online Resource~1 runs the finite check.

With this in hand the P1--P4 check becomes explicit. P1 holds because each
pairwise context supports a local assignment of the episode label.
P2 holds because overlap assignments are related by transition maps, and the
cycle product prevents global compatibility.
P3 holds because the induced context-groupoid action has no global fixed
section.
Equivalently, transporting a local label around the cycle returns its
opposite.
P4 holds relative to the following declared revision envelope:
admissible revisions may add time resolution, duplicate pairwise probes,
replace labels by conservative aliases, or add pairwise behavioral readouts,
but may not add a direct all-context token reader.
Each admissible revision commutes with the transition maps and preserves the
cycle product $-1$.
Thus no revision inside $\Omega$ dissolves the obstruction.

What would count as a genuine defeat of the obstruction also shows up here. A
probe that reads the latent episode token jointly across $P$, $W$, and $R$
changes the site. Such a probe is a non-conservative revision, and $\Omega$ must
declare it. Where it is physically and operationally admissible, the case moves
from $Z = 1$ to $Z = 0$, or to a new site requiring a fresh computation.
Otherwise the low-concern verdict stays unsupported.

This example generalizes. P1--P4 are no mere labels
pinned on after the fact. They are a checklist: context site, target fiber,
transition maps, holonomy, and revision envelope. They also pry access apart
from phenomenal attribution. The architecture can show pairwise broadcast and
report alignment while still failing to identify the global target. That is the
form of the test the later theory sections apply to IIT, GWT, and HOT at a higher
level of abstraction.

A second executable check trades the pairwise-context cycle for a small Boolean
recurrent architecture. The system has an input bit, an observable report bit, a
memory bit with update rule $m_{t+1}=x_t$, and a latent bit with update rule
$a_{t+1}=a_t$. Across every input sequence up to the checked horizon, the
report-only site sees identical traces for two states that differ only in the
latent bit. A direct latent reader resolves the target, but only by changing the
site. Online Resource~1 supplies the check as
\texttt{boolean\_recurrent\_obstruction\_check.py}.

%%─────────────────────────────────────────────────────────────────────────────
\section{A worked disclosure dossier: Llama 3.1 8B-Instruct}
\label{sec:llama-dossier}

This dossier is a positive control: it illustrates how the framework
applies to a public system without claiming LLM consciousness or its
absence. It accepts ordinary access/report evidence for a narrow target
while declining to carry that evidence over to phenomenal consciousness,
and it applies the disclosure requirement to a public system in place of a
toy architecture. The system is Meta Llama~3.1~8B-Instruct,
an open-weight, text-only,
instruction-tuned transformer released in July 2024
\citep{grattafiori2024llama3,meta2024llama31card}.
The model card describes it as an autoregressive language model using an
optimized transformer architecture, with supervised fine-tuning and RLHF for
helpfulness and safety alignment.
For the 8B model, the public record states a 128k context length, grouped-query
attention, multilingual text input, text and code output, pretraining on more
than 15 trillion tokens, and a December 2023 pretraining data cutoff.

The dossier is deliberately narrow. It concerns one self-hosted inference
instance of the released 8B instruction model, with no external tools, memory,
autonomous action loops, fine-tuning, or agent scaffolding. Any of those
additions would change the site and call for a fresh dossier.

Mechanistic interpretability already offers many ways to test whether
representations in a transformer track prompt content and shape later
output. The dossier's contribution is the audit distinction. Evidence that the model uses a context
item is evidence for an access/report target. On its own that evidence
leaves the phenomenal target untouched, and the two targets stay apart when a verdict issues.

\emph{1. Observation site.}
The admissible site contains: public model-card information; released model
weights and architecture; token-level input-output behavior; prompt
interventions; logit traces; activation probes; ablations of attention heads,
MLP blocks, layers, residual streams, and safety wrappers; and reproducible
self-hosted inference logs.
It excludes non-public training traces, vendor-internal human-feedback records,
private safety-classifier decisions, and unlogged deployment-layer state.
It also excludes biological measurements, because the system is not biological.

\emph{2. Target fiber.}
Access and phenomenal status are separate targets. For access or report, the target fiber holds
operational labels such as ``information is available for token generation,''
``the model reports a state,'' or ``the model can use a context item in later
output,'' and ordinary behavioral and mechanistic probes can test these. For
phenomenal consciousness or moral patienthood, the target fiber holds a
phenomenal-status label for an inference state, with the public site supplying no
decidable equality test for it.

\emph{3. Revision envelope.}
The declared $\Omega$ permits longer prompt batteries, adversarial prompts,
activation patching, sparse-feature probes, causal ablations, alternative
decoding parameters, quantization checks, and removal or replacement of
deployment guardrails.
It also permits comparison across self-hosted replicas of the same released
weights. It forbids adding persistent memory, tool use, online learning,
multi-agent scaffolding, biological sensors, or private training data while
still calling the result the same site. Such changes are non-conservative and
require a new $Z$ evaluation.

\emph{4. Classifier.}
For access/report targets, the classifier is behavioral and mechanistic: does a
probed representation causally affect later token generation or report behavior
under held-out prompts? For phenomenal consciousness, the public record supplies
no accepted classifier. Self-reports about consciousness are themselves outputs
under the same observation site, never labels for the hidden target.

\emph{5. Challenge record.}
A sufficient public challenge record would include prompt suites, activation
probe protocols, ablation results, negative controls, model/version hashes,
sampling parameters, deployment-wrapper descriptions, and failed attempts to
find a non-conservative probe that changes the verdict. The current public
record holds useful model-card and benchmark information, but as a
consciousness-attribution challenge log it falls short.

\emph{6. Revocation rule.}
A verdict must be recomputed once the site changes, as it does under tool
integration, persistent memory, online adaptation, fine-tuning, changed safety
wrappers, multimodal inputs, autonomous action loops, or access to private
training traces. It changes too if a future theory supplies a decidable
phenomenal classifier over this class of systems.

\emph{7. Provenance record.}
A low-concern certificate would oblige the audit to name who selected the probe
classes, who ran the red-team process, who maintained the classifier, which
infrastructure was used, and whether those actors or facilities carried
claimant-originated dependencies. Here the provenance record stays
intentionally limited: the author selected and ran the probe packet, the GPU
infrastructure was rented, the model came from an ungated public mirror after
the official gated repository proved inaccessible, and no independent red-team or
source-class provenance certificate was produced. This suffices for a worked
access/report illustration. For a low-concern phenomenal-consciousness
certificate it does not.

\emph{Verdict.}
For the access/report task in this dossier, the result is evaluable and yields
$Z=0$: one can probe, ablate, and compare the public model across replicates for
that target. For phenomenal consciousness or moral patienthood, the dossier is
non-applicable. The ground for that is the public site's lack of a target
fiber with decidable equality, a phenomenal classifier, a public
consciousness challenge log, and the provenance certificate a low-concern
verdict requires. It is no demonstration that Llama~3.1~8B-Instruct is
non-conscious. So the framework blocks a certified low-concern claim for the
phenomenal target under this dossier while letting ordinary scientific claims
about access, report, and mechanism through.

The dossier infers no
consciousness from opacity. It separates targets, treats the system as
experimentally tractable for access and report, and refuses only the stronger
phenomenal verdict, on the ground that the dossier does not contain the objects
such a verdict requires.
The executable access/report probe packet is
\texttt{llama31\_activation\_probe.py} in Online Resource~1.
It freezes the prompt schema and declared site in a manifest, then on a GPU
collects selected hidden states, reports leave-one-out probe accuracy by layer,
and runs a targeted residual-direction intervention.
This packet bears on access/report evaluability and nothing more. On 2026-05-06
it ran on an A10 GPU against the ungated
\texttt{NousResearch/Meta-Llama-3.1-8B-Instruct} mirror, using 24 controlled
key-value prompts. Token-level answer accuracy was $24/24$, with a mean
correct-label logit margin of $6.107$ against a no-mapping control margin of
roughly zero ($-0.026$). Leave-one-out centroid probes on selected hidden states
reached $1.00$ accuracy at layers 16, 24, and 31, and $0.75$ at layer 8, and a
targeted residual-direction intervention reduced the correct-label margin at
layers 16, 24, and 31 with mean deltas $-0.732$, $-1.669$, and $-1.107$. These
results support the narrow access/report claim under the declared site while
supplying no phenomenal-consciousness classifier. Online Resource~1 holds the
saved packet as \texttt{llama31\_probe\_packet\_nous/summary.json}.

The experiment functions as a positive control. Linear probes and residual
interventions are familiar tools for showing that transformer representations
track prompt content and shape later output, so the experiment adds no new
mechanistic-interpretability result. Its point lies in the audit
discipline. The framework
accepts the ordinary evidence for the access/report target, assigns it the
expected evaluable verdict, and then declines to transfer that verdict to the
phenomenal target until someone supplies the missing target fiber, classifier,
challenge record, revision envelope, and provenance certificate. Had the
access/report control failed, the dossier would have told us nothing even about
the narrow target. Because it succeeds, the non-applicability verdict for the
phenomenal target cannot be charged to a blanket skepticism about empirical AI
systems.

%%─────────────────────────────────────────────────────────────────────────────
\section{Implications for IIT, GWT, and higher-order theories}
\label{sec:implications}

These three cases do different jobs. IIT is the hardest case for the
experimentability witness, since its central quantity is defined at a level that
current measurement cannot reach. GWT and higher-order theories sit more
comfortably as theories of access and metacognition, and their difficulty
surfaces only when those narrower targets are pressed into service as evidence
for phenomenal consciousness.

\subsection{Integrated information theory}
\label{sec:iit}

IIT identifies consciousness with integrated information $\Phi$ above a threshold
\citep{tononi2016}. The live question is what one would have to measure to make a
full $\Phi$-based identification claim auditable, not whether IIT is true. Computing $\Phi$ calls for the cause-effect repertoire of every
possible subsystem, across all spatial and temporal scales, which for a system
of $N$ binary neurons means $2^N$ perturbation experiments. No current
neuroimaging or electrophysiology protocol comes close, so under any known
experimental protocol the observation record $\mathcal{R}_Y$ lacks the
information needed to compute $\Phi$ with decidable equality.

For the purpose of full $\Phi$-based identification, IIT's target $\Theta$ falls
in the prior class, because the observation protocol required to test the target
does not exist. So it is not yet in the ``$Z = 1$, identification impossible''
class. Specifically, full cause-effect repertoires lie outside
$\mathrm{Adm}(Y)$ under any currently realizable experimental protocol.

One distinction needs marking. IIT's \emph{predictions} are falsifiable, and the
2025 adversarial study challenged one salient preregistered prediction
\citep{cogitate2025} (Cogitate consortium, $n = 256$, multimodal
fMRI/MEG/iEEG): the expected sustained posterior synchronization did not appear.
The same study also challenged GNWT, through the lack of ignition at stimulus
offset and limited prefrontal representation of some conscious dimensions. Those
results are scientifically relevant, yet they leave the identification question
used here unsettled. The witness-failure verdict turns on a different question,
whether one can compute the value of $\Phi$ from experimental data. Falsifying a
theory-derived prediction needs only that the theory entail observable
consequences, with no need to compute $\Phi$. The identification problem arises only
when no observation protocol can supply the quantity the theory defines as its
target, $\Phi$ itself.
Current approximation schemes such as $\Phi^*$, MIB, and GeoMIP read best as
proxies or restricted computations. The best implementations handle small
systems far below brain scale, none of them cited here claiming decidable
equality of $\hat\Phi$ with $\Phi$ for the full target.

\textit{The tractability objection.}
An IIT proponent may answer that $\Phi$ is computable in principle from complete
cause-effect structure, the computation being exponential in $N$ yet tractable
in principle. The experimentability witness asks for something else,
\emph{decidable equality}: a procedure that definitively answers ``does this
measured quantity equal the theory-defined target $\Phi$?'' That demand is
distinct from tractability. To satisfy the witness one needs both (a)~an
observation protocol $Y$ whose record $\mathcal{R}_Y$ contains the cause-effect
repertoire, and (b)~a decidable equality check between the measured quantity and
the theoretical $\Phi$. In practice the exponential-computation barrier blocks
(a), and no current approximation scheme satisfies~(b): each $\Phi^*$ variant is
an
acknowledged lower bound or proxy, far short of a decidable equality test against
the theoretical target. So the witness-failure verdict rests on the absence of a
decidable specification. \citet{hanson2019iit} independently established a
related substrate-sensitivity problem for IIT, namely that $\Phi$ is not
invariant under state-labeling isomorphisms, so functionally identical systems
can receive different $\Phi$ values. The descent-obstruction analysis here
supplies a substrate-neutral structural reason for that failure of invariance.
IIT thus lacks the experimentability witness, meaning decidable equality on
$\Theta$, that a full $\Phi$-based identification claim would need. A version of
IIT specifying a practical approximation $\hat\Phi$ with a decidable equality
test would clear this first hurdle, and whether $\hat\Phi$ actually tracks
$\Phi$ would then become the next empirical question.


\subsection{Global workspace theory}
\label{sec:gwt}

GWT identifies consciousness with global broadcast, information made available to
multiple cognitive subsystems at once \citep{baars1988,dehaene2011}. This target
is far more tractable than $\Phi$, since masking, report, attention, and
neural-imaging paradigms can all test broadcast. For access consciousness in
Block's sense \citep{block1995}, GWT can supply the experimentability witness.
None of this lets current GNWT predictions off the hook: the Cogitate
adversarial study challenged GNWT's offset-ignition and prefrontal-content
predictions. The narrower point holds anyway, that one can specify broadcast and
access targets in an experimentable record type even where particular GNWT
predictions fail.

Trouble begins once broadcast evidence is made to license claims about
phenomenal consciousness. Two systems with identical global broadcast patterns
may differ in phenomenal experience. Where the governance target is phenomenal
consciousness, the kernel pair fires: (broadcast $+$ phenomenal experience)
against (broadcast $+$ no phenomenal experience) produce identical records under
any broadcast-based measurement. So GWT evidence can speak strongly to access
while leaving the phenomenal target unsettled.

\subsection{Higher-order theories}
\label{sec:hot}

Higher-order theories identify consciousness with a representation of being in a
mental state \citep{rosenthal2005,carruthers2000}. Here again the first step is
feasible: metacognitive reports, confidence judgments, and prefrontal-monitoring
signals all fit a record type with decidable equality. The harder
question is whether higher-order representation identifies phenomenality.
Blindsight makes the problem vivid, with subjects showing information-guided
behavior in the absence of ordinary visual experience. Should the kernel pair
(higher-order representation $+$ phenomenal experience) against (higher-order
representation $+$ no phenomenal experience) be empirically realizable,
Lemma~\ref{thm:kernelpair} applies. HOT therefore passes the witness for a
metacognitive target, and passing it leaves phenomenal consciousness still
unidentified.

%%─────────────────────────────────────────────────────────────────────────────
\section{The upload problem and evolutionary selection}
\label{sec:applications}

Consider first the upload problem. Someone proposes a digital system $D$ as
conscious by copying the behavioral profile of a biological system $B$, so that
$O(B) \approx O(D)$. Observable copying carries over the record $E_Y(h)$ while
leaving behind the descent structure, the transition maps and holonomy of the
target fiber. A verdict would require showing
$(\mathsf{C}_{Y_D}, F_D, \rho_D) \cong (\mathsf{C}_Y, F, \rho)$. Behavioral
matching alone does not reach that far, and absent the structural isomorphism the
correct verdict for $D$ is non-evaluable.
Non-evaluability here does not deny the upload moral status. It blocks the
verdict that would license disregarding the upload's welfare, so under the
precautionary premise the possible subject is protected, not cleared.

Consider next the evolutionary argument. Natural selection acts on observable
outputs and reaches $\Theta$ only through them. Where $\Theta$ is
descent-obstructed, selection on $O$ does not entail selection on $\Theta$. So
the inference ``consciousness evolved because it confers fitness advantage''
becomes structurally invalid when $Z = 1$: selection may have favored hidden
causal structures $X$ that produce adaptive $O$, with the obstructed property
$\Theta$ riding along as a byproduct.

The same machinery reads off a substrate-neutral identification claim. A digital
system that instantiates an equivalent context category, transition maps, and
holonomy structure,
$(\mathsf{C}_{Y_D}, F_D, \rho_D) \cong (\mathsf{C}_Y, F, \rho)$, falls in the
same identification class as a biological system with that structure. Descent
equivalence across the declared sites is what matters. So the theory predicts
that substrate-isomorphic systems are identification-equivalent once the site,
target fiber, and revision envelope match.

A final caution: information flow is not holonomy. A system maximizing integrated
information, global broadcast, or any other information-theoretic quantity
satisfies P1, local contextual presentability in information terms, while it may
well fail P3, nontrivial holonomy. Information-theoretic magnitude is measurable
and identifiable. What generates the obstruction is structural holonomy over the
admissible observation site, so large information flow sits comfortably with
$Z = 0$.

%%─────────────────────────────────────────────────────────────────────────────
\section{Relation to the literature}
\label{sec:literature}

We borrow the mathematics from elsewhere. Kernel-pair identification is
standard in causal inference \citep{pearl2009}, \v{C}ech cohomology for
contextuality comes from \citet{abramsky2011}, and Grothendieck descent is
sixty-year-old algebraic geometry \citep{grothendieck1971sga}. The contribution here
brings this known machinery to bear on consciousness attribution.

Older and less tidy is the philosophical setting. \citet{nagel1974} and
\citet{levine1983} made the epistemic difficulty of subjective experience hard
to ignore, and \citet{chalmers1995} gave the difficulty its contemporary name.
\citet{dennett1991} and \citet{frankish2016} press the other way, holding that
what needs explaining may be the judgments, reports, and seeming of
phenomenality, with no further hidden property posited. That dispute stays open
here. The prior question is narrower: given a declared target, can one identify it
from the admitted observations? Block's access/phenomenal distinction
\citep{block1995} marks exactly the place where many
proposed AI tests quietly switch targets.

The approach is methodological and leaves the metaphysics aside. When the
target is access, report, metacognitive monitoring, or functional integration,
ordinary experimental and mechanistic tests may identify it. When the target is
phenomenal consciousness in the sense at issue for \citet{nagel1974},
\citet{levine1983}, and \citet{chalmers1995}, the claimant has to say how that
target enters the observation site. An illusionist or behaviorist may instead
deny that any further phenomenal target is needed, which legitimately changes the
target. A replacement target is then evaluated in turn, and the discarded one
drops out. None of this settles the dispute between realist and
deflationary theories of consciousness. It is a discipline for moving
between them: state the target, state the
site, and show whether the target is identified.

The position is epistemic structuralism about scientific attribution, with no
commitment on metaphysical structuralism about what exists. A realist may hold
that there is a fact of the matter beyond the declared observations. That
possibility stands, though the observation site does not license such a fact as a
scientific verdict. An illusionist or behaviorist may hold that no further
phenomenal residue is needed, which the analysis grants. Where the proposed
target is report, judgment, seeming, or illusion-detection, that target is the
one evaluated. The error to avoid is a careless use of Block's distinction: access
evidence can be valid evidence for access while leaving phenomenal consciousness
unidentified. Which target one should care about is left open, with one
exception, that a claimant must declare a change of target.

Philosophically, the question becomes one of identification: before treating
evidence as decisive, a consciousness verdict must hold target, site, and
admissible revision envelope fixed. Lemma~\ref{thm:kernelpair} already gives a
strong version of the point for classical indistinguishability cases, and
Block's access/phenomenal distinction explains why access evidence cannot simply
migrate to the phenomenal target. The descent machinery carries the point to
contextual and non-classical substrates where kernel-pair quotients or
\v{C}ech $H^1$ are too narrow. Its fully general reduction remains a stated
obligation, with the finite factorization cases and the conservative-invariance
results establishing the part used here. From this follows a governance
corollary carrying an explicit six-element disclosure requirement that includes
the revision envelope $\Omega$. From there we evaluate major consciousness
frameworks against the experimentability precondition.

On formal frameworks for AI consciousness, \citet{butlin2023consciousness} set
the current methodological baseline by deriving computational indicator
properties from GWT, IIT, HOT, and related theories and then assessing AI
systems against them. Carefully empirical though it is, the approach presupposes
that the indicator properties are experimentally accessible.
\citet{phillips2024metamath} restate all six IIT axioms as
universal mapping properties in the sense of category theory, including the
sheaf condition.
The present framework extends theirs from definition to testability: when the sheaf
condition fails (no global sections exist), that is exactly the
descent-obstruction criterion we use as an identification-impossibility
certificate.
\citet{kleiner2021} give a categorical formulation of IIT aimed at defining
consciousness, where our aim is to characterize when such definitions become
scientifically evaluable.
\citet{gunji2026contextuality} give a recent categorical reconstruction of
contextuality in which orthomodular structure arises from classical Boolean
contexts.
That work is adjacent to our use of contextual obstruction, while our target
is identification of a hidden property from admissible observations.
\citet{campero2025framework} provide the clearest taxonomy of
impossibility-class arguments against AI consciousness, and our result
contributes a new category, identification-level impossibility independent of
substrate implementation.
\citet{phua2025test} use artificial agents and ablations to test functional
predictions associated with GWT, IIT-adjacent markers, and HOT.
Their dissociations support the need to keep access, integration, and
metacognitive targets separate before making any phenomenal-consciousness
verdict.

Closest of the adjacent prior work is \citet{inoue2026}, who develops a
sheaf-cohomological treatment of neural integration on brain manifolds, modeling
local cognitive functions as sheaf sections and classifying integration failures
as cohomological obstructions. The framework here differs from his in two
respects. First, Grothendieck effective descent serves as the obstruction
criterion in place of
\v{C}ech $H^1$. The kernel-pair example (Lemma~\ref{thm:kernelpair}) shows
descent to be strictly more general: systems exist with $H^1 = 0$ but $D = 1$,
so an $H^1$-only diagnostic misses part of the identification-failure class.
Second, the target is the identification problem, whether one can recover the
property $\Theta$ from admissible measurements, where his is neural integration.
The
Proposition (Conservative Invariance) and the revision envelope $\Omega$ bear on
the language-stability of the identification claim, a question Inou\'{e} leaves
aside.

On governance under consciousness uncertainty, \citet{long2024welfare} and
\citet{butlin2025principles} argue that precautionary governance is warranted.
The governance corollary sharpens their recommendation: for the hidden-variable
substrate class, the uncertainty is identification-theoretic under the declared
observation regime. On the normative premise adopted here, a low confidence
score of ``likely not conscious'' falls short. \citet{wolfson2026consent} and
\citet{rost2026sentience} independently anticipate related protection and
readiness concerns through non-mathematical reasoning. \citet{birch2024}
develops a graduated precautionary framework for sentience attribution, and
Corollary~\ref{cor:governance} gives a formal condition under which the
evidential basis for low concern is absent. Where \citet{zhou2025ethics} defaults
to a ``presumption of no consciousness,'' the present theorem rejects that
default once descent obstruction holds and the governance target requires
identification evidence. Under non-effectiveness, no negative verdict follows
from the absence of admissible evidence.

%%─────────────────────────────────────────────────────────────────────────────
\section{Objections}
\label{sec:objections}

The deepest objection concerns the phenomenal target itself. The descent test
needs a target fiber with decidable equality, and for phenomenal consciousness no
one has exhibited one, so a phenomenal $Z = 1$ would be either tautological or
stipulated. The premise is granted. The conclusion does not follow. This paper
claims no phenomenal $Z = 1$ for any real system. The $Z = 1$ results hold only
where the target fiber is specifiable: the broadcast-report architecture under
global episode identity, and access or report targets with decidable equality.
Where the target is phenomenal consciousness, the public site supplies no such
fiber, the descent test does not run, and the correct output is
non-applicability, the fail-closed branch, not $Z = 1$. The claim for the
phenomenal case is therefore diagnostic, not an impossibility theorem:
consciousness science cannot at present state the phenomenal target in an
auditable form, and until it can, a low-concern verdict has nothing to discharge.
Non-identifiability is nowhere assumed in order to derive the verdict. What the
framework diagnoses is non-specifiability, and that is what blocks the verdict.

A critic may object that $\Omega$ is a knob. Set it too narrow and the
obstruction is shielded from serious future tests. Set it too broad and the
result grows hard to evaluate. This is the framework's main practical pressure
point. The answer is procedural: $Z$ is always relative to a declared $\Omega$,
and the declaration forms part of the public claim, so a low-concern verdict
carries an audit burden. The claimant must say which new observables,
interventions, probes, and structural vocabularies are admissible, and why
excluded revisions fall outside the physical or operational scope. An
under-specified $\Omega$ yields neither $Z = 1$ nor $Z = 0$. It yields
non-applicability and the same fail-closed governance result.

A behaviorist offers a different objection. If two systems are indistinguishable
under every admissible observation, the behaviorist may deny that any remaining
difference is scientifically meaningful. The objection has force at the
epistemic level: when the target is not identified, science cannot license any
target attribution, positive or negative, from the declared observations. The
governance conclusion then adds a premise, that a low-concern verdict carries
moral stakes and so requires identification evidence. The
mathematics blocks the evidential claim. The normative premise fixes the
default response.

A further objection presses on substrate. The framework is substrate-neutral only
at the level of form, and substrate still does work through the declared inputs.
For each substrate, the claimant must specify the admissible site $Y$, the target
fiber $F \to H$, and the revision envelope $\Omega$. A biological system, a
digital system, and a quantum-contextual system may instantiate one descent
pattern, but the claimant has to show the equivalence. Substrate enters through
the site and the allowed revisions. Once those inputs are fixed, the theorem says
what follows. It leaves the burden of fixing them where it lies, on the claimant.

%%─────────────────────────────────────────────────────────────────────────────
\section{Scope and limitations}
\label{sec:limits}

Several things lie outside the reach of the argument. Which systems are
conscious is one of them. The theorems characterize when science cannot
answer that question. A verdict of $Z = 1$ never marks the absence of
consciousness either, only the failure of admissible experiments to settle the
matter. Nor does the \S\ref{sec:implications} evaluation pronounce IIT, GWT, or
HOT false: it is structural, and passes no truth judgment. The governance
corollary, in turn, reaches some moral-status decisions and not every one. And a
proof sketch is just that, never a proof.

The finite deterministic and exact finite stochastic factorization
results are proved in \S\ref{sec:maintheorem}. The broader reductions remain argued, not
proved: that every admissible identification procedure across all substrate
classes factors through descent, and that the descent criterion is computed
identically across all five substrate classes in Table~\ref{tab:substrates}.

Eight claims here are mechanically verified in Lean~4 (Mathlib v4.30.0-rc2),
each axiom-clean (no \texttt{sorry}, checked against an explicit axiom allowlist
of \texttt{Classical.choice}, \texttt{Quot.sound}, and \texttt{propext}): the
finite deterministic factorization biconditional ($E : H \to R$,
image-restricted); the exact finite stochastic factorization biconditional
($K : H \to \mathrm{PMF}\,R$, image-restricted, the harder of the two
factorization results, closed as separate forward and backward lemmas); the
finite kernel-pair non-factorization lemma; the finite-alias Conservative
Invariance case; a categorical sheaf-transport statement under suitable
equivalence; the Appendix~\ref{app:suite} Case~4 $Z$-formula evaluations; the
\v{C}ech-holonomy detection of \S\ref{sec:cech} on an arbitrary finite nerve (a
closed walk with nonzero holonomy admits no global section, for abelian
coefficients, with the cyclic case as a special walk); and the
revision-stability of that obstruction under cover refinement (subdividing each
edge and redistributing its label leaves the holonomy invariant and nonzero, so
no global section exists in any such refinement). This last result verifies
revision-stability only under refinement of a fixed cover. Stability under
arbitrary change of measurement language stays argued, not proved. Online
Resource~2 supplies the Lean files, which compile with only the standard axiom
dependencies. They are verification receipts for named components, not a proof of
the full descent-factorization reduction, which remains the main open reduction.

Operationalizing $\Omega$ is the principal open problem. A claimant
must declare and audit the revision envelope $\Omega$ before $Z$ is evaluable.
For real physical and computational systems, specifying and publicly
justifying $\Omega$ is open-ended.
Practitioners may under-specify, strategically exclude plausible
non-conservative revisions, or dispute which expansions count as conservative.
This leaves the theoretical result untouched, since $Z$ is well-defined relative
to any declared $\Omega$, but it lodges the evidential burden in $\Omega$'s
completeness and public auditability. The practical next step is therefore an
$\Omega$ audit protocol: a declaration schema, an adversarial challenge process
for missing revisions, and a revocation rule for when a low-concern verdict
rests on an under-specified or strategically narrow envelope. Such a protocol
must also refuse a common shortcut. Success on seeded omissions qualifies a
challenger's competence while leaving unbounded the chance that adversarial
probe classes remain unsurfaced. Residual escape must be bounded under a declared
adversary model and public challenge process, on pain of non-applicability. The
same holds for independence. A provenance certificate asks for more than the
absence of a visible conflict: the dossier must certify the source-class
provenance of auditors, red teams, classifiers, panel members, infrastructure
providers, and benchmark or data providers, and it must aggregate
claimant-originated benefits across proxy and sub-threshold paths. Lacking that
protocol, the framework can still state non-applicability cleanly even while a
real disclosure dossier stays too easy to game.

Treat internal adversarial review as a source of candidate objections. It does
not certify the theorem. Evidence for the framework has to come from public
proofs, public code, worked systems, or external audit records. Several such
checks are now complete.
\S\ref{sec:llama-dossier} gives a public-system dossier for an open-weight
transformer and a reproducible activation-probe packet for the access/report
target;
the finite stochastic script checks the exact kernel case;
the Boolean recurrent script runs the obstruction calculation on a small
transition system; and
the prediction packet states the following falsifiable disclosure-packet
prediction: for a fixed text-only, self-hosted, open-weight transformer with no
external tools, memory, online learning, or autonomous action loop,
access/report use of an input item should be evaluable from public weights,
token-level behavior, logits, hidden states, activation probes, causal
ablations, prompt interventions, and reproducible inference logs. Phenomenal
consciousness or moral patienthood should remain non-applicable under that
same site unless the claimant supplies a target fiber with decidable equality,
a public classifier, a public challenge record, a declared revision envelope,
and a provenance certificate for the challenge apparatus. The prediction fails
if an auditor supplies those phenomenal-target objects without adding a
non-conservative probe, or if the access/report packet cannot produce
above-chance behavioral evidence for a system known to use the relevant input.

Several lines of further work remain. The public-system dossier should
be extended to a deployed system with tool use, memory, and logged deployment
state. The finite proof should be pushed from exact kernels to finite-sample
statistical decision procedures and finite measurable systems before anyone
claims the fully general version. The prediction packet should go to an external
registry or independent auditor before it is used as external evidence. And the
philosophy-of-mind engagement in \S\ref{sec:literature} should continue by
stating in each case exactly which target is at issue, whether access, report,
integration, metacognition, or phenomenal attribution.

%%─────────────────────────────────────────────────────────────────────────────
\section{Conclusion}
\label{sec:conclusion}

For some hidden-variable systems, consciousness attribution carries an
identification problem that no further observation inside the declared site can
remove. This happens when the target fiber has revision-stable non-effective
descent. The systems in this class are characterized by four structural
properties, P1--P4, that generate necessary non-identification for any substrate
supplying the required site, target fiber, descent semantics, and revision
envelope. Under the descent-factorization obligation, the revision-stable
descent criterion (Theorem~\ref{thm:descent}) states the impossibility, and the
Proposition (Conservative Invariance) shows that conservative re-presentations of
the structural language leave the obstruction in place. The framework's boundary
is also fixed: the bridge from P1--P4 to consciousness attribution calls
for an exogenous commitment or a non-conservative revision of the admissible
site. The governance corollary follows from all this. When $Z = 1$, a
low-concern verdict lacks the identification evidence it requires, and an
organization that issues such a verdict before satisfying the six-element
disclosure requirement, the declared revision envelope $\Omega$ among them, has
not supplied the evidence its verdict needs.

The framework becomes usable only once the claimant writes down the observation
site, target fiber, transition data, and revision envelope for the system
at hand. The worked application in \S\ref{sec:worked} shows the shape of that
specification on a small case. Without those objects, any consciousness verdict
is premature.

%%─────────────────────────────────────────────────────────────────────────────
\bibliographystyle{plainnat}
\bibliography{refs}

%%─────────────────────────────────────────────────────────────────────────────
\appendix

\clearpage
\section{Falsification suite}
\label{app:suite}

The code below confirms all the quantitative predictions and runs with no
dependencies beyond Python's standard library.
The suite tests four claims: (1)~the contextuality triangle has $D = 1$ and
$\check{H}^1 = 1$;
(2)~the kernel-pair case has $D = 1$ and $\check{H}^1 = 0$, demonstrating
$\check{H}^1$-only recall $= 0.5$; (3)~the identifiable control has $D = 0$;
and (4)~the revision-stability predictions, namely that a conservative
language expansion leaves $D$ unchanged and a non-conservative
hidden-state-revealing revision removes the obstruction.

\begin{lstlisting}
import itertools

def rank_f2(matrix):
    m = [row[:] for row in matrix]
    if not m:
        return 0
    rows, cols = len(m), len(m[0])
    rank = 0
    for col in range(cols):
        pivot = next(
            (r for r in range(rank, rows) if m[r][col] % 2 == 1), None)
        if pivot is None:
            continue
        m[rank], m[pivot] = m[pivot], m[rank]
        for r in range(rows):
            if r != rank and m[r][col] % 2 == 1:
                m[r] = [(m[r][c] ^ m[rank][c]) for c in range(cols)]
        rank += 1
    return rank

# Case 1: contextuality triangle
CONTEXTS = ["C12", "C23", "C13"]
CONTEXT_STATES = {
    "C12": ("h1","h2"), "C23": ("h2","h3"), "C13": ("h1","h3")}

def overlapping_states(ca, cb):
    return set(CONTEXT_STATES[ca]) & set(CONTEXT_STATES[cb])

def local_value(ctx, state, section):
    return section[CONTEXT_STATES[ctx].index(state)]

def is_globally_consistent(sections):
    for ca, cb in itertools.combinations(CONTEXTS, 2):
        for state in overlapping_states(ca, cb):
            if local_value(ca, state, sections[ca]) != \
               local_value(cb, state, sections[cb]):
                return False
    return True

ODD_PARITY = {
    ctx: [(a,b) for a in (0,1) for b in (0,1) if (a^b)==1]
    for ctx in CONTEXTS}

def global_sections(local_family):
    return [
        {ctx: sec for ctx, sec in zip(CONTEXTS, combo)}
        for combo in itertools.product(
            *[local_family[c] for c in CONTEXTS])
        if is_globally_consistent(
            {ctx: sec for ctx, sec in zip(CONTEXTS, combo)})]

coboundary    = [[1,1,0],[0,1,1],[1,0,1]]
triangle_rank = rank_f2(coboundary)           # 2
triangle_h1   = 3 - triangle_rank             # 1
triangle_gs   = len(global_sections(ODD_PARITY))  # 0
triangle_D    = 1 if triangle_gs == 0 else 0      # 1

assert triangle_rank == 2
assert triangle_h1   == 1
assert triangle_gs   == 0
assert triangle_D    == 1

# Case 2: kernel-pair (H^1 = 0, D = 1)
obs    = {"a": "same", "b": "same"}
target = {"a": 0, "b": 1}

def kernel_pair_D(obs, theta):
    for x, y in itertools.combinations(obs, 2):
        if obs[x] == obs[y] and theta[x] != theta[y]:
            return 1
    return 0

kernel_h1 = 0
kernel_D  = kernel_pair_D(obs, target)    # 1
assert kernel_h1 == 0
assert kernel_D  == 1

# Case 3: identifiable control
obs_ctrl  = {"a": "obs_a", "b": "obs_b"}
tgt_ctrl  = {"a": 0, "b": 1}
control_D = kernel_pair_D(obs_ctrl, tgt_ctrl)  # 0
assert control_D == 0

# Detector comparison
total          = triangle_D + kernel_D
h1_detected    = (1 if triangle_h1>0 else 0) + (1 if kernel_h1>0 else 0)
h1_recall      = h1_detected / total      # 0.5
descent_recall = total / total            # 1.0

assert h1_recall      == 0.5
assert descent_recall == 1.0
\end{lstlisting}

Case 4 verifies a definition and runs no independent
empirical test. By a category-theoretic argument, the Conservative
Invariance proposition guarantees that adding a Morita-equivalent alias preserves
the descent flag, while the non-conservative variant adds an observable that
distinguishes the two hidden states and so, by the definition of the descent
obstruction, must resolve it. Case 4 is included as a verification benchmark,
confirming that the formal machinery behaves as the proposition predicts. It adds no
evidence beyond what the proposition's category-theoretic content already
establishes.

\begin{lstlisting}
# Case 4: revision-stability benchmark
# Base: kernel-pair obstruction (D_base = 1).
# Conservative revision: add a definitional alias of the existing observable.
#   The Proposition (Conservative Invariance) predicts D unchanged.
# Non-conservative revision: add an observable that distinguishes the hidden states.
#   This changes Y; the obstruction is resolved (D = 0).

def compute_D(obs_map, theta_map):
    for x, y in itertools.combinations(obs_map, 2):
        if obs_map[x] == obs_map[y] and theta_map[x] != theta_map[y]:
            return 1
    return 0

obs_base   = {"a": "obs_X", "b": "obs_X"}
theta_base = {"a": 0, "b": 1}
D_base     = compute_D(obs_base, theta_base)
assert D_base == 1, "Base kernel-pair must have obstruction"

obs_conservative   = {"a": ("obs_X","obs_X_alias"), "b": ("obs_X","obs_X_alias")}
theta_conservative = theta_base
D_conservative     = compute_D(obs_conservative, theta_conservative)
assert D_conservative == 1, "Conservative expansion must preserve obstruction"
assert D_conservative - D_base == 0, "Delta D for conservative revision must be 0"

obs_nonconservative   = {"a": ("obs_X","reveals_a"), "b": ("obs_X","reveals_b")}
theta_nonconservative = theta_base
D_nonconservative     = compute_D(obs_nonconservative, theta_nonconservative)
assert D_nonconservative == 0, "Non-conservative revision must remove obstruction"
assert D_nonconservative - D_base == -1, "Delta D for non-conservative revision must be -1"

conservative_overturns    = 0
nonconservative_overturns = 1
actual_conservative_overturns    = 1 - D_conservative
actual_nonconservative_overturns = 1 - D_nonconservative
assert actual_conservative_overturns    == conservative_overturns
assert actual_nonconservative_overturns == nonconservative_overturns

n_generators     = 10
m_hidden_states  = 100
revision_subsets = 2 ** n_generators
hidden_pairs     = m_hidden_states * (m_hidden_states - 1) // 2
total_checks     = revision_subsets * hidden_pairs
assert revision_subsets == 1024
assert hidden_pairs     == 4950
assert total_checks     == 5_068_800

if __name__ == "__main__":
    print("All assertions passed.")
    print(f"Triangle: rank={triangle_rank}, H1={triangle_h1}, "
          f"global_sections={triangle_gs}, D={triangle_D}")
    print(f"Kernel-pair: H1={kernel_h1}, D={kernel_D}")
    print(f"H1-only recall={h1_recall} | Descent recall={descent_recall}")
    print(f"D_base={D_base}, D_conservative={D_conservative}, "
          f"D_nonconservative={D_nonconservative}")
    print(f"Conservative overturns={actual_conservative_overturns}, "
          f"Non-conservative overturns={actual_nonconservative_overturns}")
    print(f"Revision audit cost (n={n_generators}, m={m_hidden_states}): "
          f"{total_checks:,} pair-checks")
\end{lstlisting}

\end{document}
